\scanpage{107}
% L1--5
\uline{Corollaire} 2 \deleted{Sous les conditions du corollaire 1,} Supposons $X$ \deleted{de plus} quasi compact \inserted{et (T$_0$)}, et que \deleted{tte partie} $X$ ait une base formée d'ouverts \deleted{\unsure{qu. cpts|qu.-cpts}} rétrocompacts, qui sont qu. cpts pour la topologie constructible. Soit $E$ une partie de $X$. Conditions équivalentes

% L6--7
a) $E$ est fermé pour la top. constructible [i.e. est intersection \deleted{pour la} de parties constructibles]

% L8--11
b) $E$ est compact pour la top. constructible, i.e. pour tte famille de parties constructibles $(E_i)_{i\in I}$ de $X$, telle que $E\subset\bigcup_{i\in I}E_i$, il existe une sous-famille finie $J\subset I$, telle que $E\subset\bigcup_{i\in J}E_i$.

\margin{-- \unsure{généralt|général}, \uline{partie \unsure{antitypique|antypique|atypique|analytique}} $=$ \uline{intersection de loclt constructibles} Nature locale de la notion ??}
% slanted

% L12
\deleted{$X$ \unsure{est}} En effet, $X^{\mathrm{cons}}$ est compact, donc a) $\Longleftrightarrow$ b).

% L13--16
On dit parfois \inserted{alors} que \deleted{$X$} $E$ est une \uline{partie antitypique de} $X$. On peut donner un \unsure{énoncé} analogue si $X$ n'est plus supposé qu cpt [mais satisfaisant aux autres hyp]

% L17--20
\[
\begin{aligned}
& E\ \text{fermé \inserted{dans $X^{\mathrm{cons}}$} [i.e. intersection de parties loc. constructibles]}\\
& \Updownarrow\\
& E\ \text{loc cpt pour la top. constructible i.e. pour tt}\\
& \text{\deleted{partie \unsure{ouverte}} }U_i,\quad E\cap \deleted{\unsure{E|X}}U_i\ \text{satisfait la condition b}\\
& \text{ci-dessus dans } U_i .
\end{aligned}
\]

% L21--25
\uline{Corollaire} 3 Sous les conditions du corollaire 2, Conditions équivalentes

a) $E$ et $X-E$ sont antitypiques

b) $E$ est \deleted{constructible} ouvert et fermé dans $X^{\mathrm{cons}}$

c) $E$ est constructible.
% archive sheet number: 106
